\section*{Chapitre \ref{ch:g11:dna-replication} --- La réplication de l'ADN}

\begin{solution}{exo:g11:dna-replication:1}
Chaque brin sert de matrice sur laquelle un nouveau brin complémentaire
est bâti. Cela ne fonctionne que parce que A s'apparie avec T et G avec
C~: la séquence d'un brin fixe celle de l'autre, si bien qu'un brin
seul porte toute l'information.
\end{solution}

\begin{solution}{exo:g11:dna-replication:2}
Semi-conservative~: chaque molécule fille garde un brin parental et
reçoit un brin neuf. Conservative~: la molécule parentale reste entière
et un double brin entièrement neuf est bâti à côté d'elle.
\end{solution}

\begin{solution}{exo:g11:dna-replication:3}
La phase S du cycle cellulaire~; la quantité d'\omterm{def:g10:universal-dna:information}{ADN} par \omterm{def:g10:cells-common-unit:organelle}{noyau} double.
\end{solution}

\begin{solution}{exo:g11:dna-replication:4}
L'\omterm{prop:g11:dna-replication:forks}{ADN polymérase} ajoute à un brin en croissance les \omterm{def:g10:universal-dna:nucleotide}{nucléotides}
complémentaires de la matrice, un à la fois. Une \omterm{prop:g11:dna-replication:forks}{fourche de réplication}
est le point mobile où les brins parentaux sont séparés et où les deux
brins neufs sont en cours de construction.
\end{solution}

\begin{solution}{exo:g11:dna-replication:5}
\texttt{CTAATGTCCG}, base sous base.
\end{solution}

\begin{solution}{exo:g11:dna-replication:6}
Conservatif, génération 1~: une bande lourde et une bande légère, pas
d'hybride --- or seule une bande hybride a été vue. Dispersif,
génération 2~: une bande unique au quart du chemin du léger vers le
lourd --- or deux bandes distinctes ont été vues. Chaque modèle prévoit
des bandes qui n'ont pas été observées.
\end{solution}

\begin{solution}{exo:g11:dna-replication:7}
$2^5 = 32$ molécules, 2 hybrides et 30 légères~: une bande légère
quinze fois la bande hybride, aucune bande lourde.
\end{solution}

\begin{solution}{exo:g11:dna-replication:8}
Génération 0~: une bande légère. Génération 1~: une bande hybride.
Génération 2~: bandes hybride et lourde d'égale intensité --- l'image
en miroir de l'expérience d'origine.
\end{solution}

\begin{solution}{exo:g11:dna-replication:9}
$\num{2.5e8}/50 = \num{5e6}\,\unit{s}$, soit environ 58 jours --- bien
plus que 8 heures. Le \omterm{def:g10:universal-dna:information}{chromosome} doit être copié à partir de nombreuses
origines à la fois, chacune ouvrant deux fourches.
\end{solution}

\begin{solution}{exo:g11:dna-replication:10}
Les deux brins de $\num{3.2e9}$ paires de \omterm{def:g10:universal-dna:nucleotide}{bases}, deux fois~:
$\num{6.4e9}$ paires de \omterm{def:g10:universal-dna:nucleotide}{bases}, soit $\num{6.4e9}$ \omterm{def:g10:universal-dna:nucleotide}{nucléotides} neufs
enchaînés (un par paire de \omterm{def:g10:universal-dna:nucleotide}{bases} copiée). À $10^{-9}$ par \omterm{def:g10:universal-dna:nucleotide}{nucléotide},
environ 6 erreurs.
\end{solution}

\begin{solution}{exo:g11:dna-replication:11}
Ouvrir un œil expose de la matrice des deux côtés de l'origine~; chaque
fourche copie en s'éloignant de l'origine, si bien que les deux vont en
sens opposés et couvrent ensemble toute la région. Quand deux yeux se
rencontrent, leurs fourches fusionnent et les brins fils se raboutent
en molécules continues.
\end{solution}

\begin{solution}{exo:g11:dna-replication:12}
Les bactéries cessent de se diviser (pas de réplication, pas de
division). La plaie cesse de cicatriser, puisque sa réparation demande
des divisions cellulaires. Le neurone n'est pas touché~: il ne réplique
plus son \omterm{def:g10:universal-dna:information}{ADN}. Un médicament qui arrête les \omterm{def:g10:cells-common-unit:cell}{cellules} en division rapide
tout en épargnant celles qui ne se divisent pas suggère un traitement
contre les bactéries ou contre les \omterm{def:g10:cells-common-unit:cell}{cellules} cancéreuses en division.
\end{solution}

\begin{solution}{exo:g11:dna-replication:13}
Les deux brins d'origine ne sont jamais détruits~: à chaque génération,
chacun est de nouveau apparié à un brin léger neuf, si bien qu'il
existe toujours exactement deux molécules hybrides. Après 10
générations~: $2/2^{10} = 2/1024
\approx 0.2\%$ --- présentes, mais trop pâles pour être vues.
\end{solution}

\begin{solution}{exo:g11:dna-replication:14}
$\num{6.4e9} \times 10^{-5} = \num{64000}$ erreurs par division. Avec
vingt mille \omterm{def:g10:universal-dna:gene}{gènes}, des centaines de \omterm{def:g10:universal-dna:gene}{gènes} seraient endommagés à chaque
division~; après quelques divisions, les \omterm{def:g10:cells-common-unit:cell}{cellules} porteraient des
défauts létaux et la lignée s'éteindrait.
\end{solution}

\begin{solution}{exo:g11:dna-replication:15}
Un nouveau tour de réplication démarre à l'origine avant que le
précédent ne soit terminé~: le \omterm{def:g10:universal-dna:information}{chromosome} porte à la fois plusieurs
paires de fourches emboîtées. Chaque division reçoit alors un
\omterm{def:g10:universal-dna:information}{chromosome} déjà partiellement copié pour la suivante, et les divisions
peuvent se succéder plus vite qu'une copie complète ne dure.
\end{solution}

\begin{solution}{pb:g11:dna-replication:1}
\textbf{1.} L'azote fait 16~\% de la masse et est 7~\% plus lourd~:
$0.16 \times 0.07 \approx 1.1\%$. La méthode du gradient sépare des
densités qui diffèrent d'une fraction de pour cent, et 1~\% donne donc
des bandes bien séparées.

\textbf{2.} Pour que la quasi-totalité de l'\omterm{def:g10:universal-dna:information}{ADN}, et non la moitié ou
les trois quarts, soit lourde~: après 14 générations, il reste moins
d'une part sur $10^4$ des brins légers d'origine.

\textbf{3.} Exactement à mi-chemin entre les deux bandes~: sa densité
est la moyenne.

\textbf{4.} Une seule bactérie détient quelques femtogrammes d'\omterm{def:g10:universal-dna:information}{ADN}~; la
photographie en demande des microgrammes, donc des milliards de
\omterm{def:g10:cells-common-unit:cell}{cellules}.

\textbf{5.} En gardant une culture en milieu lourd comme témoin et en
vérifiant que sa bande restait lourde, et en vérifiant la bande légère
de bactéries cultivées en milieu léger d'un bout à l'autre.

\textbf{6.} Semi-conservatif~: génération 1, une bande hybride~;
génération 2, bandes hybride et légère, égales.

\textbf{7.} Conservatif~: génération 1, bandes lourde et légère,
égales~; génération 2, une lourde pour trois légères.

\textbf{8.} Dispersif~: génération 1, une bande à la position
hybride~; génération 2, une bande au quart du chemin du léger au lourd.

\textbf{9.} Le modèle conservatif, qui ne prévoyait aucune bande
hybride, est éliminé. Le semi-conservatif et le dispersif prévoient
tous deux la bande hybride unique et restent en lice.

\textbf{10.} Le semi-conservatif survit~: lui seul prévoit deux bandes
séparées. Le modèle dispersif prévoit une bande unique qui s'allège
sans jamais se dédoubler.

\textbf{11.} $2^n$ molécules, dont exactement 2 portent un brin
d'origine.

\textbf{12.} Léger : hybride $= (2^n - 2) : 2$~: génération 3, $3:1$~;
génération 4, $7:1$~; génération 6, $31:1$.

\textbf{13.} La fraction hybride $2/2^n < 0.05$ dès que $2^n > 40$~:
à partir de la génération 6 ($2/64 \approx 3\%$).

\textbf{14.} Deux bandes, une lourde et une légère, en quantités
égales~: une molécule hybride est un brin lourd plus un brin léger.
Cela met à l'épreuve le modèle dispersif, qui prévoit que chaque brin
isolé est de densité intermédiaire, et l'élimine indépendamment.

\textbf{15.} En partant du léger, une génération en milieu lourd donne
une bande hybride et deux générations donnent hybride plus lourde~: le
même schéma en miroir, ce qui montre que le résultat ne dépend pas de
l'isotope qui joue le rôle de l'«~ancien~».

\textbf{16.} $\num{4.6e6}/(2 \times 1000) = \qty{2300}{s}$, soit
environ 38 minutes.

\textbf{17.} 77 minutes. Sur un cercle, les deux brins sont exposés des
deux côtés de toute ouverture~; deux fourches quittant une origine en
sens opposés est la façon la plus simple de couvrir tout le cercle.

\textbf{18.} Une origine copie $2 \times 50 \times 8 \times 3600 =
\num{2.88e6}$ paires de \omterm{def:g10:universal-dna:nucleotide}{bases} en 8 heures~; $\num{6.4e9}/\num{2.88e6}
\approx 2200$ origines au strict minimum (en réalité des dizaines de
milliers, qui s'activent à des moments différents).

\textbf{19.} Bactérie~: $\num{9.2e6} \times 10^{-9} \approx 0.01$
erreur par réplication~; \omterm{def:g10:cells-common-unit:cell}{cellule} humaine~: environ 6. L'essentiel du
génome humain n'est pas fait de \omterm{def:g10:universal-dna:gene}{gènes}, si bien que la plupart des
erreurs tombent là où elles ne changent rien, et la \omterm{def:g10:cells-common-unit:cell}{cellule} possède
deux exemplaires de chaque \omterm{def:g10:universal-dna:gene}{gène}.

\textbf{20.} Les deux bandes de la deuxième génération --- hybride et
légère en quantités égales --- ont prouvé que la réplication est
semi-conservative~; les deux fourches copient le \omterm{def:g10:universal-dna:information}{chromosome} bactérien
en environ 38 minutes, à peu près le temps de génération de la
bactérie, de sorte que la copie règle le rythme de la division.
\end{solution}
